\documentclass[11pt]{article}

    \usepackage[breakable]{tcolorbox}
    \usepackage{parskip} % Stop auto-indenting (to mimic markdown behaviour)
    

    % Basic figure setup, for now with no caption control since it's done
    % automatically by Pandoc (which extracts ![](path) syntax from Markdown).
    \usepackage{graphicx}
    % Keep aspect ratio if custom image width or height is specified
    \setkeys{Gin}{keepaspectratio}
    % Maintain compatibility with old templates. Remove in nbconvert 6.0
    \let\Oldincludegraphics\includegraphics
    % Ensure that by default, figures have no caption (until we provide a
    % proper Figure object with a Caption API and a way to capture that
    % in the conversion process - todo).
    \usepackage{caption}
    \DeclareCaptionFormat{nocaption}{}
    \captionsetup{format=nocaption,aboveskip=0pt,belowskip=0pt}

    \usepackage{float}
    \floatplacement{figure}{H} % forces figures to be placed at the correct location
    \usepackage{xcolor} % Allow colors to be defined
    \usepackage{enumerate} % Needed for markdown enumerations to work
    \usepackage{geometry} % Used to adjust the document margins
    \usepackage{amsmath} % Equations
    \usepackage{amssymb} % Equations
    \usepackage{textcomp} % defines textquotesingle
    % Hack from http://tex.stackexchange.com/a/47451/13684:
    \AtBeginDocument{%
        \def\PYZsq{\textquotesingle}% Upright quotes in Pygmentized code
    }
    \usepackage{upquote} % Upright quotes for verbatim code
    \usepackage{eurosym} % defines \euro

    \usepackage{iftex}
    \ifPDFTeX
        \usepackage[T1]{fontenc}
        \IfFileExists{alphabeta.sty}{
              \usepackage{alphabeta}
          }{
              \usepackage[mathletters]{ucs}
              \usepackage[utf8x]{inputenc}
          }
    \else
        \usepackage{fontspec}
        \usepackage{unicode-math}
    \fi

    \usepackage{fancyvrb} % verbatim replacement that allows latex
    \usepackage{grffile} % extends the file name processing of package graphics
                         % to support a larger range
    \makeatletter % fix for old versions of grffile with XeLaTeX
    \@ifpackagelater{grffile}{2019/11/01}
    {
      % Do nothing on new versions
    }
    {
      \def\Gread@@xetex#1{%
        \IfFileExists{"\Gin@base".bb}%
        {\Gread@eps{\Gin@base.bb}}%
        {\Gread@@xetex@aux#1}%
      }
    }
    \makeatother
    \usepackage[Export]{adjustbox} % Used to constrain images to a maximum size
    \adjustboxset{max size={0.9\linewidth}{0.9\paperheight}}

    % The hyperref package gives us a pdf with properly built
    % internal navigation ('pdf bookmarks' for the table of contents,
    % internal cross-reference links, web links for URLs, etc.)
    \usepackage{hyperref}
    % The default LaTeX title has an obnoxious amount of whitespace. By default,
    % titling removes some of it. It also provides customization options.
    \usepackage{titling}
    \usepackage{longtable} % longtable support required by pandoc >1.10
    \usepackage{booktabs}  % table support for pandoc > 1.12.2
    \usepackage{array}     % table support for pandoc >= 2.11.3
    \usepackage{calc}      % table minipage width calculation for pandoc >= 2.11.1
    \usepackage[inline]{enumitem} % IRkernel/repr support (it uses the enumerate* environment)
    \usepackage[normalem]{ulem} % ulem is needed to support strikethroughs (\sout)
                                % normalem makes italics be italics, not underlines
    \usepackage{soul}      % strikethrough (\st) support for pandoc >= 3.0.0
    \usepackage{mathrsfs}
    

    
    % Colors for the hyperref package
    \definecolor{urlcolor}{rgb}{0,.145,.698}
    \definecolor{linkcolor}{rgb}{.71,0.21,0.01}
    \definecolor{citecolor}{rgb}{.12,.54,.11}

    % ANSI colors
    \definecolor{ansi-black}{HTML}{3E424D}
    \definecolor{ansi-black-intense}{HTML}{282C36}
    \definecolor{ansi-red}{HTML}{E75C58}
    \definecolor{ansi-red-intense}{HTML}{B22B31}
    \definecolor{ansi-green}{HTML}{00A250}
    \definecolor{ansi-green-intense}{HTML}{007427}
    \definecolor{ansi-yellow}{HTML}{DDB62B}
    \definecolor{ansi-yellow-intense}{HTML}{B27D12}
    \definecolor{ansi-blue}{HTML}{208FFB}
    \definecolor{ansi-blue-intense}{HTML}{0065CA}
    \definecolor{ansi-magenta}{HTML}{D160C4}
    \definecolor{ansi-magenta-intense}{HTML}{A03196}
    \definecolor{ansi-cyan}{HTML}{60C6C8}
    \definecolor{ansi-cyan-intense}{HTML}{258F8F}
    \definecolor{ansi-white}{HTML}{C5C1B4}
    \definecolor{ansi-white-intense}{HTML}{A1A6B2}
    \definecolor{ansi-default-inverse-fg}{HTML}{FFFFFF}
    \definecolor{ansi-default-inverse-bg}{HTML}{000000}

    % common color for the border for error outputs.
    \definecolor{outerrorbackground}{HTML}{FFDFDF}

    % commands and environments needed by pandoc snippets
    % extracted from the output of `pandoc -s`
    \providecommand{\tightlist}{%
      \setlength{\itemsep}{0pt}\setlength{\parskip}{0pt}}
    \DefineVerbatimEnvironment{Highlighting}{Verbatim}{commandchars=\\\{\}}
    % Add ',fontsize=\small' for more characters per line
    \newenvironment{Shaded}{}{}
    \newcommand{\KeywordTok}[1]{\textcolor[rgb]{0.00,0.44,0.13}{\textbf{{#1}}}}
    \newcommand{\DataTypeTok}[1]{\textcolor[rgb]{0.56,0.13,0.00}{{#1}}}
    \newcommand{\DecValTok}[1]{\textcolor[rgb]{0.25,0.63,0.44}{{#1}}}
    \newcommand{\BaseNTok}[1]{\textcolor[rgb]{0.25,0.63,0.44}{{#1}}}
    \newcommand{\FloatTok}[1]{\textcolor[rgb]{0.25,0.63,0.44}{{#1}}}
    \newcommand{\CharTok}[1]{\textcolor[rgb]{0.25,0.44,0.63}{{#1}}}
    \newcommand{\StringTok}[1]{\textcolor[rgb]{0.25,0.44,0.63}{{#1}}}
    \newcommand{\CommentTok}[1]{\textcolor[rgb]{0.38,0.63,0.69}{\textit{{#1}}}}
    \newcommand{\OtherTok}[1]{\textcolor[rgb]{0.00,0.44,0.13}{{#1}}}
    \newcommand{\AlertTok}[1]{\textcolor[rgb]{1.00,0.00,0.00}{\textbf{{#1}}}}
    \newcommand{\FunctionTok}[1]{\textcolor[rgb]{0.02,0.16,0.49}{{#1}}}
    \newcommand{\RegionMarkerTok}[1]{{#1}}
    \newcommand{\ErrorTok}[1]{\textcolor[rgb]{1.00,0.00,0.00}{\textbf{{#1}}}}
    \newcommand{\NormalTok}[1]{{#1}}

    % Additional commands for more recent versions of Pandoc
    \newcommand{\ConstantTok}[1]{\textcolor[rgb]{0.53,0.00,0.00}{{#1}}}
    \newcommand{\SpecialCharTok}[1]{\textcolor[rgb]{0.25,0.44,0.63}{{#1}}}
    \newcommand{\VerbatimStringTok}[1]{\textcolor[rgb]{0.25,0.44,0.63}{{#1}}}
    \newcommand{\SpecialStringTok}[1]{\textcolor[rgb]{0.73,0.40,0.53}{{#1}}}
    \newcommand{\ImportTok}[1]{{#1}}
    \newcommand{\DocumentationTok}[1]{\textcolor[rgb]{0.73,0.13,0.13}{\textit{{#1}}}}
    \newcommand{\AnnotationTok}[1]{\textcolor[rgb]{0.38,0.63,0.69}{\textbf{\textit{{#1}}}}}
    \newcommand{\CommentVarTok}[1]{\textcolor[rgb]{0.38,0.63,0.69}{\textbf{\textit{{#1}}}}}
    \newcommand{\VariableTok}[1]{\textcolor[rgb]{0.10,0.09,0.49}{{#1}}}
    \newcommand{\ControlFlowTok}[1]{\textcolor[rgb]{0.00,0.44,0.13}{\textbf{{#1}}}}
    \newcommand{\OperatorTok}[1]{\textcolor[rgb]{0.40,0.40,0.40}{{#1}}}
    \newcommand{\BuiltInTok}[1]{{#1}}
    \newcommand{\ExtensionTok}[1]{{#1}}
    \newcommand{\PreprocessorTok}[1]{\textcolor[rgb]{0.74,0.48,0.00}{{#1}}}
    \newcommand{\AttributeTok}[1]{\textcolor[rgb]{0.49,0.56,0.16}{{#1}}}
    \newcommand{\InformationTok}[1]{\textcolor[rgb]{0.38,0.63,0.69}{\textbf{\textit{{#1}}}}}
    \newcommand{\WarningTok}[1]{\textcolor[rgb]{0.38,0.63,0.69}{\textbf{\textit{{#1}}}}}
    \makeatletter
    \newsavebox\pandoc@box
    \newcommand*\pandocbounded[1]{%
      \sbox\pandoc@box{#1}%
      % scaling factors for width and height
      \Gscale@div\@tempa\textheight{\dimexpr\ht\pandoc@box+\dp\pandoc@box\relax}%
      \Gscale@div\@tempb\linewidth{\wd\pandoc@box}%
      % select the smaller of both
      \ifdim\@tempb\p@<\@tempa\p@
        \let\@tempa\@tempb
      \fi
      % scaling accordingly (\@tempa < 1)
      \ifdim\@tempa\p@<\p@
        \scalebox{\@tempa}{\usebox\pandoc@box}%
      % scaling not needed, use as it is
      \else
        \usebox{\pandoc@box}%
      \fi
    }
    \makeatother

    % Define a nice break command that doesn't care if a line doesn't already
    % exist.
    \def\br{\hspace*{\fill} \\* }
    % Math Jax compatibility definitions
    \def\gt{>}
    \def\lt{<}
    \let\Oldtex\TeX
    \let\Oldlatex\LaTeX
    \renewcommand{\TeX}{\textrm{\Oldtex}}
    \renewcommand{\LaTeX}{\textrm{\Oldlatex}}
    % Document parameters
    % Document title
    \title{ps3p2}
    
    
    
    
    
    
    
% Pygments definitions
\makeatletter
\def\PY@reset{\let\PY@it=\relax \let\PY@bf=\relax%
    \let\PY@ul=\relax \let\PY@tc=\relax%
    \let\PY@bc=\relax \let\PY@ff=\relax}
\def\PY@tok#1{\csname PY@tok@#1\endcsname}
\def\PY@toks#1+{\ifx\relax#1\empty\else%
    \PY@tok{#1}\expandafter\PY@toks\fi}
\def\PY@do#1{\PY@bc{\PY@tc{\PY@ul{%
    \PY@it{\PY@bf{\PY@ff{#1}}}}}}}
\def\PY#1#2{\PY@reset\PY@toks#1+\relax+\PY@do{#2}}

\@namedef{PY@tok@w}{\def\PY@tc##1{\textcolor[rgb]{0.73,0.73,0.73}{##1}}}
\@namedef{PY@tok@c}{\let\PY@it=\textit\def\PY@tc##1{\textcolor[rgb]{0.24,0.48,0.48}{##1}}}
\@namedef{PY@tok@cp}{\def\PY@tc##1{\textcolor[rgb]{0.61,0.40,0.00}{##1}}}
\@namedef{PY@tok@k}{\let\PY@bf=\textbf\def\PY@tc##1{\textcolor[rgb]{0.00,0.50,0.00}{##1}}}
\@namedef{PY@tok@kp}{\def\PY@tc##1{\textcolor[rgb]{0.00,0.50,0.00}{##1}}}
\@namedef{PY@tok@kt}{\def\PY@tc##1{\textcolor[rgb]{0.69,0.00,0.25}{##1}}}
\@namedef{PY@tok@o}{\def\PY@tc##1{\textcolor[rgb]{0.40,0.40,0.40}{##1}}}
\@namedef{PY@tok@ow}{\let\PY@bf=\textbf\def\PY@tc##1{\textcolor[rgb]{0.67,0.13,1.00}{##1}}}
\@namedef{PY@tok@nb}{\def\PY@tc##1{\textcolor[rgb]{0.00,0.50,0.00}{##1}}}
\@namedef{PY@tok@nf}{\def\PY@tc##1{\textcolor[rgb]{0.00,0.00,1.00}{##1}}}
\@namedef{PY@tok@nc}{\let\PY@bf=\textbf\def\PY@tc##1{\textcolor[rgb]{0.00,0.00,1.00}{##1}}}
\@namedef{PY@tok@nn}{\let\PY@bf=\textbf\def\PY@tc##1{\textcolor[rgb]{0.00,0.00,1.00}{##1}}}
\@namedef{PY@tok@ne}{\let\PY@bf=\textbf\def\PY@tc##1{\textcolor[rgb]{0.80,0.25,0.22}{##1}}}
\@namedef{PY@tok@nv}{\def\PY@tc##1{\textcolor[rgb]{0.10,0.09,0.49}{##1}}}
\@namedef{PY@tok@no}{\def\PY@tc##1{\textcolor[rgb]{0.53,0.00,0.00}{##1}}}
\@namedef{PY@tok@nl}{\def\PY@tc##1{\textcolor[rgb]{0.46,0.46,0.00}{##1}}}
\@namedef{PY@tok@ni}{\let\PY@bf=\textbf\def\PY@tc##1{\textcolor[rgb]{0.44,0.44,0.44}{##1}}}
\@namedef{PY@tok@na}{\def\PY@tc##1{\textcolor[rgb]{0.41,0.47,0.13}{##1}}}
\@namedef{PY@tok@nt}{\let\PY@bf=\textbf\def\PY@tc##1{\textcolor[rgb]{0.00,0.50,0.00}{##1}}}
\@namedef{PY@tok@nd}{\def\PY@tc##1{\textcolor[rgb]{0.67,0.13,1.00}{##1}}}
\@namedef{PY@tok@s}{\def\PY@tc##1{\textcolor[rgb]{0.73,0.13,0.13}{##1}}}
\@namedef{PY@tok@sd}{\let\PY@it=\textit\def\PY@tc##1{\textcolor[rgb]{0.73,0.13,0.13}{##1}}}
\@namedef{PY@tok@si}{\let\PY@bf=\textbf\def\PY@tc##1{\textcolor[rgb]{0.64,0.35,0.47}{##1}}}
\@namedef{PY@tok@se}{\let\PY@bf=\textbf\def\PY@tc##1{\textcolor[rgb]{0.67,0.36,0.12}{##1}}}
\@namedef{PY@tok@sr}{\def\PY@tc##1{\textcolor[rgb]{0.64,0.35,0.47}{##1}}}
\@namedef{PY@tok@ss}{\def\PY@tc##1{\textcolor[rgb]{0.10,0.09,0.49}{##1}}}
\@namedef{PY@tok@sx}{\def\PY@tc##1{\textcolor[rgb]{0.00,0.50,0.00}{##1}}}
\@namedef{PY@tok@m}{\def\PY@tc##1{\textcolor[rgb]{0.40,0.40,0.40}{##1}}}
\@namedef{PY@tok@gh}{\let\PY@bf=\textbf\def\PY@tc##1{\textcolor[rgb]{0.00,0.00,0.50}{##1}}}
\@namedef{PY@tok@gu}{\let\PY@bf=\textbf\def\PY@tc##1{\textcolor[rgb]{0.50,0.00,0.50}{##1}}}
\@namedef{PY@tok@gd}{\def\PY@tc##1{\textcolor[rgb]{0.63,0.00,0.00}{##1}}}
\@namedef{PY@tok@gi}{\def\PY@tc##1{\textcolor[rgb]{0.00,0.52,0.00}{##1}}}
\@namedef{PY@tok@gr}{\def\PY@tc##1{\textcolor[rgb]{0.89,0.00,0.00}{##1}}}
\@namedef{PY@tok@ge}{\let\PY@it=\textit}
\@namedef{PY@tok@gs}{\let\PY@bf=\textbf}
\@namedef{PY@tok@ges}{\let\PY@bf=\textbf\let\PY@it=\textit}
\@namedef{PY@tok@gp}{\let\PY@bf=\textbf\def\PY@tc##1{\textcolor[rgb]{0.00,0.00,0.50}{##1}}}
\@namedef{PY@tok@go}{\def\PY@tc##1{\textcolor[rgb]{0.44,0.44,0.44}{##1}}}
\@namedef{PY@tok@gt}{\def\PY@tc##1{\textcolor[rgb]{0.00,0.27,0.87}{##1}}}
\@namedef{PY@tok@err}{\def\PY@bc##1{{\setlength{\fboxsep}{\string -\fboxrule}\fcolorbox[rgb]{1.00,0.00,0.00}{1,1,1}{\strut ##1}}}}
\@namedef{PY@tok@kc}{\let\PY@bf=\textbf\def\PY@tc##1{\textcolor[rgb]{0.00,0.50,0.00}{##1}}}
\@namedef{PY@tok@kd}{\let\PY@bf=\textbf\def\PY@tc##1{\textcolor[rgb]{0.00,0.50,0.00}{##1}}}
\@namedef{PY@tok@kn}{\let\PY@bf=\textbf\def\PY@tc##1{\textcolor[rgb]{0.00,0.50,0.00}{##1}}}
\@namedef{PY@tok@kr}{\let\PY@bf=\textbf\def\PY@tc##1{\textcolor[rgb]{0.00,0.50,0.00}{##1}}}
\@namedef{PY@tok@bp}{\def\PY@tc##1{\textcolor[rgb]{0.00,0.50,0.00}{##1}}}
\@namedef{PY@tok@fm}{\def\PY@tc##1{\textcolor[rgb]{0.00,0.00,1.00}{##1}}}
\@namedef{PY@tok@vc}{\def\PY@tc##1{\textcolor[rgb]{0.10,0.09,0.49}{##1}}}
\@namedef{PY@tok@vg}{\def\PY@tc##1{\textcolor[rgb]{0.10,0.09,0.49}{##1}}}
\@namedef{PY@tok@vi}{\def\PY@tc##1{\textcolor[rgb]{0.10,0.09,0.49}{##1}}}
\@namedef{PY@tok@vm}{\def\PY@tc##1{\textcolor[rgb]{0.10,0.09,0.49}{##1}}}
\@namedef{PY@tok@sa}{\def\PY@tc##1{\textcolor[rgb]{0.73,0.13,0.13}{##1}}}
\@namedef{PY@tok@sb}{\def\PY@tc##1{\textcolor[rgb]{0.73,0.13,0.13}{##1}}}
\@namedef{PY@tok@sc}{\def\PY@tc##1{\textcolor[rgb]{0.73,0.13,0.13}{##1}}}
\@namedef{PY@tok@dl}{\def\PY@tc##1{\textcolor[rgb]{0.73,0.13,0.13}{##1}}}
\@namedef{PY@tok@s2}{\def\PY@tc##1{\textcolor[rgb]{0.73,0.13,0.13}{##1}}}
\@namedef{PY@tok@sh}{\def\PY@tc##1{\textcolor[rgb]{0.73,0.13,0.13}{##1}}}
\@namedef{PY@tok@s1}{\def\PY@tc##1{\textcolor[rgb]{0.73,0.13,0.13}{##1}}}
\@namedef{PY@tok@mb}{\def\PY@tc##1{\textcolor[rgb]{0.40,0.40,0.40}{##1}}}
\@namedef{PY@tok@mf}{\def\PY@tc##1{\textcolor[rgb]{0.40,0.40,0.40}{##1}}}
\@namedef{PY@tok@mh}{\def\PY@tc##1{\textcolor[rgb]{0.40,0.40,0.40}{##1}}}
\@namedef{PY@tok@mi}{\def\PY@tc##1{\textcolor[rgb]{0.40,0.40,0.40}{##1}}}
\@namedef{PY@tok@il}{\def\PY@tc##1{\textcolor[rgb]{0.40,0.40,0.40}{##1}}}
\@namedef{PY@tok@mo}{\def\PY@tc##1{\textcolor[rgb]{0.40,0.40,0.40}{##1}}}
\@namedef{PY@tok@ch}{\let\PY@it=\textit\def\PY@tc##1{\textcolor[rgb]{0.24,0.48,0.48}{##1}}}
\@namedef{PY@tok@cm}{\let\PY@it=\textit\def\PY@tc##1{\textcolor[rgb]{0.24,0.48,0.48}{##1}}}
\@namedef{PY@tok@cpf}{\let\PY@it=\textit\def\PY@tc##1{\textcolor[rgb]{0.24,0.48,0.48}{##1}}}
\@namedef{PY@tok@c1}{\let\PY@it=\textit\def\PY@tc##1{\textcolor[rgb]{0.24,0.48,0.48}{##1}}}
\@namedef{PY@tok@cs}{\let\PY@it=\textit\def\PY@tc##1{\textcolor[rgb]{0.24,0.48,0.48}{##1}}}

\def\PYZbs{\char`\\}
\def\PYZus{\char`\_}
\def\PYZob{\char`\{}
\def\PYZcb{\char`\}}
\def\PYZca{\char`\^}
\def\PYZam{\char`\&}
\def\PYZlt{\char`\<}
\def\PYZgt{\char`\>}
\def\PYZsh{\char`\#}
\def\PYZpc{\char`\%}
\def\PYZdl{\char`\$}
\def\PYZhy{\char`\-}
\def\PYZsq{\char`\'}
\def\PYZdq{\char`\"}
\def\PYZti{\char`\~}
% for compatibility with earlier versions
\def\PYZat{@}
\def\PYZlb{[}
\def\PYZrb{]}
\makeatother


    % For linebreaks inside Verbatim environment from package fancyvrb.
    \makeatletter
        \newbox\Wrappedcontinuationbox
        \newbox\Wrappedvisiblespacebox
        \newcommand*\Wrappedvisiblespace {\textcolor{red}{\textvisiblespace}}
        \newcommand*\Wrappedcontinuationsymbol {\textcolor{red}{\llap{\tiny$\m@th\hookrightarrow$}}}
        \newcommand*\Wrappedcontinuationindent {3ex }
        \newcommand*\Wrappedafterbreak {\kern\Wrappedcontinuationindent\copy\Wrappedcontinuationbox}
        % Take advantage of the already applied Pygments mark-up to insert
        % potential linebreaks for TeX processing.
        %        {, <, #, %, $, ' and ": go to next line.
        %        _, }, ^, &, >, - and ~: stay at end of broken line.
        % Use of \textquotesingle for straight quote.
        \newcommand*\Wrappedbreaksatspecials {%
            \def\PYGZus{\discretionary{\char`\_}{\Wrappedafterbreak}{\char`\_}}%
            \def\PYGZob{\discretionary{}{\Wrappedafterbreak\char`\{}{\char`\{}}%
            \def\PYGZcb{\discretionary{\char`\}}{\Wrappedafterbreak}{\char`\}}}%
            \def\PYGZca{\discretionary{\char`\^}{\Wrappedafterbreak}{\char`\^}}%
            \def\PYGZam{\discretionary{\char`\&}{\Wrappedafterbreak}{\char`\&}}%
            \def\PYGZlt{\discretionary{}{\Wrappedafterbreak\char`\<}{\char`\<}}%
            \def\PYGZgt{\discretionary{\char`\>}{\Wrappedafterbreak}{\char`\>}}%
            \def\PYGZsh{\discretionary{}{\Wrappedafterbreak\char`\#}{\char`\#}}%
            \def\PYGZpc{\discretionary{}{\Wrappedafterbreak\char`\%}{\char`\%}}%
            \def\PYGZdl{\discretionary{}{\Wrappedafterbreak\char`\$}{\char`\$}}%
            \def\PYGZhy{\discretionary{\char`\-}{\Wrappedafterbreak}{\char`\-}}%
            \def\PYGZsq{\discretionary{}{\Wrappedafterbreak\textquotesingle}{\textquotesingle}}%
            \def\PYGZdq{\discretionary{}{\Wrappedafterbreak\char`\"}{\char`\"}}%
            \def\PYGZti{\discretionary{\char`\~}{\Wrappedafterbreak}{\char`\~}}%
        }
        % Some characters . , ; ? ! / are not pygmentized.
        % This macro makes them "active" and they will insert potential linebreaks
        \newcommand*\Wrappedbreaksatpunct {%
            \lccode`\~`\.\lowercase{\def~}{\discretionary{\hbox{\char`\.}}{\Wrappedafterbreak}{\hbox{\char`\.}}}%
            \lccode`\~`\,\lowercase{\def~}{\discretionary{\hbox{\char`\,}}{\Wrappedafterbreak}{\hbox{\char`\,}}}%
            \lccode`\~`\;\lowercase{\def~}{\discretionary{\hbox{\char`\;}}{\Wrappedafterbreak}{\hbox{\char`\;}}}%
            \lccode`\~`\:\lowercase{\def~}{\discretionary{\hbox{\char`\:}}{\Wrappedafterbreak}{\hbox{\char`\:}}}%
            \lccode`\~`\?\lowercase{\def~}{\discretionary{\hbox{\char`\?}}{\Wrappedafterbreak}{\hbox{\char`\?}}}%
            \lccode`\~`\!\lowercase{\def~}{\discretionary{\hbox{\char`\!}}{\Wrappedafterbreak}{\hbox{\char`\!}}}%
            \lccode`\~`\/\lowercase{\def~}{\discretionary{\hbox{\char`\/}}{\Wrappedafterbreak}{\hbox{\char`\/}}}%
            \catcode`\.\active
            \catcode`\,\active
            \catcode`\;\active
            \catcode`\:\active
            \catcode`\?\active
            \catcode`\!\active
            \catcode`\/\active
            \lccode`\~`\~
        }
    \makeatother

    \let\OriginalVerbatim=\Verbatim
    \makeatletter
    \renewcommand{\Verbatim}[1][1]{%
        %\parskip\z@skip
        \sbox\Wrappedcontinuationbox {\Wrappedcontinuationsymbol}%
        \sbox\Wrappedvisiblespacebox {\FV@SetupFont\Wrappedvisiblespace}%
        \def\FancyVerbFormatLine ##1{\hsize\linewidth
            \vtop{\raggedright\hyphenpenalty\z@\exhyphenpenalty\z@
                \doublehyphendemerits\z@\finalhyphendemerits\z@
                \strut ##1\strut}%
        }%
        % If the linebreak is at a space, the latter will be displayed as visible
        % space at end of first line, and a continuation symbol starts next line.
        % Stretch/shrink are however usually zero for typewriter font.
        \def\FV@Space {%
            \nobreak\hskip\z@ plus\fontdimen3\font minus\fontdimen4\font
            \discretionary{\copy\Wrappedvisiblespacebox}{\Wrappedafterbreak}
            {\kern\fontdimen2\font}%
        }%

        % Allow breaks at special characters using \PYG... macros.
        \Wrappedbreaksatspecials
        % Breaks at punctuation characters . , ; ? ! and / need catcode=\active
        \OriginalVerbatim[#1,codes*=\Wrappedbreaksatpunct]%
    }
    \makeatother

    % Exact colors from NB
    \definecolor{incolor}{HTML}{303F9F}
    \definecolor{outcolor}{HTML}{D84315}
    \definecolor{cellborder}{HTML}{CFCFCF}
    \definecolor{cellbackground}{HTML}{F7F7F7}

    % prompt
    \makeatletter
    \newcommand{\boxspacing}{\kern\kvtcb@left@rule\kern\kvtcb@boxsep}
    \makeatother
    \newcommand{\prompt}[4]{
        {\ttfamily\llap{{\color{#2}[#3]:\hspace{3pt}#4}}\vspace{-\baselineskip}}
    }
    

    
    % Prevent overflowing lines due to hard-to-break entities
    \sloppy
    % Setup hyperref package
    \hypersetup{
      breaklinks=true,  % so long urls are correctly broken across lines
      colorlinks=true,
      urlcolor=urlcolor,
      linkcolor=linkcolor,
      citecolor=citecolor,
      }
    % Slightly bigger margins than the latex defaults
    
    \geometry{verbose,tmargin=1in,bmargin=1in,lmargin=1in,rmargin=1in}
    
    

\begin{document}
    
    \maketitle
    
    

    
    Consider the model studied in Problem 3 on Problem set 2, where
\((Y, X)\) follow the logit
model\[Y_1 = \mathbf{1}\{\beta_{01} + X_1\beta_{02} + V_1 \geq 0\}, \quad (1)\]where
\(V_1\) is an unobservable variable independent of \(X_1\) with the
logit c.d.f. \(\Lambda(v) = \frac{\exp(v)}{1+\exp(v)}\), and \(X_1\) is
a Bernoulli random variable such that \(P(X_1 = 1) = 0.75\). The
parameter \(\beta_0 = (\beta_{01}, \beta_{02})\) denotes the true
population values and \(\beta = (\beta_1, \beta_2)\) denotes a generic
parameter vector, possibly different from \(\beta_0\). Suppose
\(\beta_{01} = \beta_{02} = 1\).Write a function named LogitDraws that
draws 1000 realizations of \(V_1\) and \(X_1\). For each draw of \(V_1\)
and \(X_1\) then compute \(Y_1\) according to (1), and output from the
function two \(1000 \times 1\) vectors \(y\) and \(x\) containing all
the realizations of \(Y_1\) and \(X_1\), respectively.

    \begin{tcolorbox}[breakable, size=fbox, boxrule=1pt, pad at break*=1mm,colback=cellbackground, colframe=cellborder]
\prompt{In}{incolor}{1}{\boxspacing}
\begin{Verbatim}[commandchars=\\\{\}]
\PY{k+kn}{import}\PY{+w}{ }\PY{n+nn}{numpy}\PY{+w}{ }\PY{k}{as}\PY{+w}{ }\PY{n+nn}{np}

\PY{k}{def}\PY{+w}{ }\PY{n+nf}{LogitDraws}\PY{p}{(}\PY{p}{)}\PY{p}{:}
\PY{+w}{    }\PY{l+s+sd}{\PYZdq{}\PYZdq{}\PYZdq{}}
\PY{l+s+sd}{    Draws 1000 realizations of V\PYZus{}1 and X\PYZus{}1, computes Y\PYZus{}1, and returns the draws.}

\PY{l+s+sd}{    The model is Y\PYZus{}1 = 1\PYZob{}beta\PYZus{}01 + X\PYZus{}1*beta\PYZus{}02 + V\PYZus{}1 \PYZgt{}= 0\PYZcb{},}
\PY{l+s+sd}{    where V\PYZus{}1 \PYZti{} Logistic(0, 1), X\PYZus{}1 \PYZti{} Bernoulli(0.75),}
\PY{l+s+sd}{    and beta\PYZus{}01 = beta\PYZus{}02 = 1.}

\PY{l+s+sd}{    Returns:}
\PY{l+s+sd}{        y (np.ndarray): 1000x1 vector of Y\PYZus{}1 realizations.}
\PY{l+s+sd}{        x (np.ndarray): 1000x1 vector of X\PYZus{}1 realizations.}
\PY{l+s+sd}{    \PYZdq{}\PYZdq{}\PYZdq{}}
    \PY{c+c1}{\PYZsh{} Number of draws}
    \PY{n}{n\PYZus{}draws} \PY{o}{=} \PY{l+m+mi}{1000}

    \PY{c+c1}{\PYZsh{} True parameter values}
    \PY{n}{beta\PYZus{}01} \PY{o}{=} \PY{l+m+mi}{1}
    \PY{n}{beta\PYZus{}02} \PY{o}{=} \PY{l+m+mi}{1}

    \PY{c+c1}{\PYZsh{} Draw 1000 realizations of X\PYZus{}1 from a Bernoulli distribution with p=0.75}
    \PY{n}{x} \PY{o}{=} \PY{n}{np}\PY{o}{.}\PY{n}{random}\PY{o}{.}\PY{n}{binomial}\PY{p}{(}\PY{l+m+mi}{1}\PY{p}{,} \PY{l+m+mf}{0.75}\PY{p}{,} \PY{n}{n\PYZus{}draws}\PY{p}{)}

    \PY{c+c1}{\PYZsh{} Draw 1000 realizations of V\PYZus{}1 from a standard logistic distribution}
    \PY{n}{v} \PY{o}{=} \PY{n}{np}\PY{o}{.}\PY{n}{random}\PY{o}{.}\PY{n}{logistic}\PY{p}{(}\PY{n}{loc}\PY{o}{=}\PY{l+m+mi}{0}\PY{p}{,} \PY{n}{scale}\PY{o}{=}\PY{l+m+mi}{1}\PY{p}{,} \PY{n}{size}\PY{o}{=}\PY{n}{n\PYZus{}draws}\PY{p}{)}

    \PY{c+c1}{\PYZsh{} Compute Y\PYZus{}1}
    \PY{n}{y} \PY{o}{=} \PY{p}{(}\PY{n}{beta\PYZus{}01} \PY{o}{+} \PY{n}{x} \PY{o}{*} \PY{n}{beta\PYZus{}02} \PY{o}{+} \PY{n}{v} \PY{o}{\PYZgt{}}\PY{o}{=} \PY{l+m+mi}{0}\PY{p}{)}\PY{o}{.}\PY{n}{astype}\PY{p}{(}\PY{n+nb}{int}\PY{p}{)}

    \PY{c+c1}{\PYZsh{} Reshape to be column vectors (1000x1)}
    \PY{n}{y} \PY{o}{=} \PY{n}{y}\PY{o}{.}\PY{n}{reshape}\PY{p}{(}\PY{o}{\PYZhy{}}\PY{l+m+mi}{1}\PY{p}{,} \PY{l+m+mi}{1}\PY{p}{)}
    \PY{n}{x} \PY{o}{=} \PY{n}{x}\PY{o}{.}\PY{n}{reshape}\PY{p}{(}\PY{o}{\PYZhy{}}\PY{l+m+mi}{1}\PY{p}{,} \PY{l+m+mi}{1}\PY{p}{)}

    \PY{k}{return} \PY{n}{y}\PY{p}{,} \PY{n}{x}
\end{Verbatim}
\end{tcolorbox}

    \begin{tcolorbox}[breakable, size=fbox, boxrule=1pt, pad at break*=1mm,colback=cellbackground, colframe=cellborder]
\prompt{In}{incolor}{2}{\boxspacing}
\begin{Verbatim}[commandchars=\\\{\}]
\PY{c+c1}{\PYZsh{} Call the function and store the results}
\PY{n}{y\PYZus{}draws}\PY{p}{,} \PY{n}{x\PYZus{}draws} \PY{o}{=} \PY{n}{LogitDraws}\PY{p}{(}\PY{p}{)}

\PY{c+c1}{\PYZsh{} Print the shapes to verify}
\PY{n+nb}{print}\PY{p}{(}\PY{l+s+s2}{\PYZdq{}}\PY{l+s+s2}{Shape of y\PYZus{}draws:}\PY{l+s+s2}{\PYZdq{}}\PY{p}{,} \PY{n}{y\PYZus{}draws}\PY{o}{.}\PY{n}{shape}\PY{p}{)}
\PY{n+nb}{print}\PY{p}{(}\PY{l+s+s2}{\PYZdq{}}\PY{l+s+s2}{Shape of x\PYZus{}draws:}\PY{l+s+s2}{\PYZdq{}}\PY{p}{,} \PY{n}{x\PYZus{}draws}\PY{o}{.}\PY{n}{shape}\PY{p}{)}

\PY{c+c1}{\PYZsh{} Print the first 10 draws as a sample}
\PY{n+nb}{print}\PY{p}{(}\PY{l+s+s2}{\PYZdq{}}\PY{l+s+se}{\PYZbs{}n}\PY{l+s+s2}{Sample draws (first 10):}\PY{l+s+s2}{\PYZdq{}}\PY{p}{)}
\PY{n+nb}{print}\PY{p}{(}\PY{l+s+s2}{\PYZdq{}}\PY{l+s+s2}{Y | X}\PY{l+s+s2}{\PYZdq{}}\PY{p}{)}
\PY{n+nb}{print}\PY{p}{(}\PY{l+s+s2}{\PYZdq{}}\PY{l+s+s2}{\PYZhy{}\PYZhy{}\PYZhy{}\PYZhy{}\PYZhy{}}\PY{l+s+s2}{\PYZdq{}}\PY{p}{)}
\PY{k}{for} \PY{n}{i} \PY{o+ow}{in} \PY{n+nb}{range}\PY{p}{(}\PY{l+m+mi}{10}\PY{p}{)}\PY{p}{:}
    \PY{n+nb}{print}\PY{p}{(}\PY{l+s+sa}{f}\PY{l+s+s2}{\PYZdq{}}\PY{l+s+si}{\PYZob{}}\PY{n}{y\PYZus{}draws}\PY{p}{[}\PY{n}{i}\PY{p}{]}\PY{p}{[}\PY{l+m+mi}{0}\PY{p}{]}\PY{l+s+si}{\PYZcb{}}\PY{l+s+s2}{ | }\PY{l+s+si}{\PYZob{}}\PY{n}{x\PYZus{}draws}\PY{p}{[}\PY{n}{i}\PY{p}{]}\PY{p}{[}\PY{l+m+mi}{0}\PY{p}{]}\PY{l+s+si}{\PYZcb{}}\PY{l+s+s2}{\PYZdq{}}\PY{p}{)}
\end{Verbatim}
\end{tcolorbox}

    \begin{Verbatim}[commandchars=\\\{\}]
Shape of y\_draws: (1000, 1)
Shape of x\_draws: (1000, 1)

Sample draws (first 10):
Y | X
-----
1 | 1
1 | 0
1 | 1
1 | 1
0 | 1
1 | 1
1 | 0
1 | 1
1 | 0
1 | 0
    \end{Verbatim}

    \begin{tcolorbox}[breakable, size=fbox, boxrule=1pt, pad at break*=1mm,colback=cellbackground, colframe=cellborder]
\prompt{In}{incolor}{3}{\boxspacing}
\begin{Verbatim}[commandchars=\\\{\}]
\PY{k+kn}{import}\PY{+w}{ }\PY{n+nn}{scipy}\PY{n+nn}{.}\PY{n+nn}{optimize}\PY{+w}{ }\PY{k}{as}\PY{+w}{ }\PY{n+nn}{opt}
\PY{k+kn}{from}\PY{+w}{ }\PY{n+nn}{scipy}\PY{+w}{ }\PY{k+kn}{import} \PY{n}{special} \PY{c+c1}{\PYZsh{} Use special.expit for the logistic function}

\PY{c+c1}{\PYZsh{} Assuming y\PYZus{}draws and x\PYZus{}draws are available from the previous cell}

\PY{c+c1}{\PYZsh{} Add a constant to the exogenous variables for the intercept}
\PY{n}{X} \PY{o}{=} \PY{n}{np}\PY{o}{.}\PY{n}{hstack}\PY{p}{(}\PY{p}{(}\PY{n}{np}\PY{o}{.}\PY{n}{ones}\PY{p}{(}\PY{p}{(}\PY{n}{x\PYZus{}draws}\PY{o}{.}\PY{n}{shape}\PY{p}{[}\PY{l+m+mi}{0}\PY{p}{]}\PY{p}{,} \PY{l+m+mi}{1}\PY{p}{)}\PY{p}{)}\PY{p}{,} \PY{n}{x\PYZus{}draws}\PY{p}{)}\PY{p}{)}
\PY{n}{y} \PY{o}{=} \PY{n}{y\PYZus{}draws}\PY{o}{.}\PY{n}{flatten}\PY{p}{(}\PY{p}{)}

\PY{c+c1}{\PYZsh{} Define the negative log\PYZhy{}likelihood function for Logit}
\PY{k}{def}\PY{+w}{ }\PY{n+nf}{neg\PYZus{}log\PYZus{}likelihood\PYZus{}logit}\PY{p}{(}\PY{n}{beta}\PY{p}{,} \PY{n}{y}\PY{p}{,} \PY{n}{X}\PY{p}{)}\PY{p}{:}
    \PY{c+c1}{\PYZsh{} Calculate the linear combination of predictors}
    \PY{n}{linear\PYZus{}pred} \PY{o}{=} \PY{n}{X} \PY{o}{@} \PY{n}{beta}
    \PY{c+c1}{\PYZsh{} Calculate the probability of Y=1 using the logistic CDF (also called expit or sigmoid)}
    \PY{n}{p\PYZus{}y1} \PY{o}{=} \PY{n}{special}\PY{o}{.}\PY{n}{expit}\PY{p}{(}\PY{n}{linear\PYZus{}pred}\PY{p}{)}
    \PY{c+c1}{\PYZsh{} The log\PYZhy{}likelihood function}
    \PY{c+c1}{\PYZsh{} To avoid log(0) errors, we use a small epsilon}
    \PY{n}{epsilon} \PY{o}{=} \PY{l+m+mf}{1e\PYZhy{}9}
    \PY{n}{log\PYZus{}likelihood} \PY{o}{=} \PY{n}{np}\PY{o}{.}\PY{n}{sum}\PY{p}{(}\PY{n}{y} \PY{o}{*} \PY{n}{np}\PY{o}{.}\PY{n}{log}\PY{p}{(}\PY{n}{p\PYZus{}y1} \PY{o}{+} \PY{n}{epsilon}\PY{p}{)} \PY{o}{+} \PY{p}{(}\PY{l+m+mi}{1} \PY{o}{\PYZhy{}} \PY{n}{y}\PY{p}{)} \PY{o}{*} \PY{n}{np}\PY{o}{.}\PY{n}{log}\PY{p}{(}\PY{l+m+mi}{1} \PY{o}{\PYZhy{}} \PY{n}{p\PYZus{}y1} \PY{o}{+} \PY{n}{epsilon}\PY{p}{)}\PY{p}{)}
    \PY{c+c1}{\PYZsh{} Return the negative of the log\PYZhy{}likelihood because we are minimizing}
    \PY{k}{return} \PY{o}{\PYZhy{}}\PY{n}{log\PYZus{}likelihood}

\PY{c+c1}{\PYZsh{} Initial guess for the parameters (beta\PYZus{}1, beta\PYZus{}2)}
\PY{n}{initial\PYZus{}beta} \PY{o}{=} \PY{n}{np}\PY{o}{.}\PY{n}{zeros}\PY{p}{(}\PY{n}{X}\PY{o}{.}\PY{n}{shape}\PY{p}{[}\PY{l+m+mi}{1}\PY{p}{]}\PY{p}{)}

\PY{c+c1}{\PYZsh{} Perform the MLE estimation}
\PY{n}{result} \PY{o}{=} \PY{n}{opt}\PY{o}{.}\PY{n}{minimize}\PY{p}{(}
    \PY{n}{neg\PYZus{}log\PYZus{}likelihood\PYZus{}logit}\PY{p}{,} \PY{c+c1}{\PYZsh{} Use the correct logit function}
    \PY{n}{initial\PYZus{}beta}\PY{p}{,}
    \PY{n}{args}\PY{o}{=}\PY{p}{(}\PY{n}{y}\PY{p}{,} \PY{n}{X}\PY{p}{)}\PY{p}{,}
    \PY{n}{method}\PY{o}{=}\PY{l+s+s1}{\PYZsq{}}\PY{l+s+s1}{BFGS}\PY{l+s+s1}{\PYZsq{}}
\PY{p}{)}

\PY{k}{if} \PY{n}{result}\PY{o}{.}\PY{n}{success}\PY{p}{:}
    \PY{n}{mle\PYZus{}beta} \PY{o}{=} \PY{n}{result}\PY{o}{.}\PY{n}{x}
    \PY{n+nb}{print}\PY{p}{(}\PY{l+s+s2}{\PYZdq{}}\PY{l+s+s2}{MLE estimates for beta (Logit Model):}\PY{l+s+s2}{\PYZdq{}}\PY{p}{)}
    \PY{n+nb}{print}\PY{p}{(}\PY{l+s+sa}{f}\PY{l+s+s2}{\PYZdq{}}\PY{l+s+s2}{beta\PYZus{}1 (intercept): }\PY{l+s+si}{\PYZob{}}\PY{n}{mle\PYZus{}beta}\PY{p}{[}\PY{l+m+mi}{0}\PY{p}{]}\PY{l+s+si}{:}\PY{l+s+s2}{.4f}\PY{l+s+si}{\PYZcb{}}\PY{l+s+s2}{\PYZdq{}}\PY{p}{)}
    \PY{n+nb}{print}\PY{p}{(}\PY{l+s+sa}{f}\PY{l+s+s2}{\PYZdq{}}\PY{l+s+s2}{beta\PYZus{}2 (coefficient for X\PYZus{}1): }\PY{l+s+si}{\PYZob{}}\PY{n}{mle\PYZus{}beta}\PY{p}{[}\PY{l+m+mi}{1}\PY{p}{]}\PY{l+s+si}{:}\PY{l+s+s2}{.4f}\PY{l+s+si}{\PYZcb{}}\PY{l+s+s2}{\PYZdq{}}\PY{p}{)}
\PY{k}{else}\PY{p}{:}
    \PY{n+nb}{print}\PY{p}{(}\PY{l+s+s2}{\PYZdq{}}\PY{l+s+s2}{MLE estimation failed.}\PY{l+s+s2}{\PYZdq{}}\PY{p}{)}
    \PY{n+nb}{print}\PY{p}{(}\PY{n}{result}\PY{o}{.}\PY{n}{message}\PY{p}{)}
\end{Verbatim}
\end{tcolorbox}

    \begin{Verbatim}[commandchars=\\\{\}]
MLE estimates for beta (Logit Model):
beta\_1 (intercept): 0.8457
beta\_2 (coefficient for X\_1): 0.8895
    \end{Verbatim}

    \begin{tcolorbox}[breakable, size=fbox, boxrule=1pt, pad at break*=1mm,colback=cellbackground, colframe=cellborder]
\prompt{In}{incolor}{4}{\boxspacing}
\begin{Verbatim}[commandchars=\\\{\}]
\PY{c+c1}{\PYZsh{} This cell calculates the asymptotic variance of the MLE estimators using the BHHH estimator.}
\PY{c+c1}{\PYZsh{} It depends on the \PYZsq{}mle\PYZus{}beta\PYZsq{}, \PYZsq{}y\PYZsq{}, and \PYZsq{}X\PYZsq{} variables from the previous cell.}

\PY{k+kn}{import}\PY{+w}{ }\PY{n+nn}{numpy}\PY{+w}{ }\PY{k}{as}\PY{+w}{ }\PY{n+nn}{np}
\PY{k+kn}{from}\PY{+w}{ }\PY{n+nn}{scipy}\PY{+w}{ }\PY{k+kn}{import} \PY{n}{special}

\PY{c+c1}{\PYZsh{} Calculate p\PYZus{}i for each observation using the estimated beta}
\PY{n}{z} \PY{o}{=} \PY{n}{X} \PY{o}{@} \PY{n}{mle\PYZus{}beta}
\PY{n}{p} \PY{o}{=} \PY{n}{special}\PY{o}{.}\PY{n}{expit}\PY{p}{(}\PY{n}{z}\PY{p}{)} \PY{c+c1}{\PYZsh{} Use logistic CDF}

\PY{c+c1}{\PYZsh{} To avoid division by zero in case p is 0 or 1}
\PY{n}{epsilon} \PY{o}{=} \PY{l+m+mf}{1e\PYZhy{}9}
\PY{n}{p} \PY{o}{=} \PY{n}{np}\PY{o}{.}\PY{n}{clip}\PY{p}{(}\PY{n}{p}\PY{p}{,} \PY{n}{epsilon}\PY{p}{,} \PY{l+m+mi}{1} \PY{o}{\PYZhy{}} \PY{n}{epsilon}\PY{p}{)}

\PY{c+c1}{\PYZsh{} Calculate the score (gradient) for each observation for the Logit model}
\PY{c+c1}{\PYZsh{} g is a (n\PYZus{}obs, n\PYZus{}params) matrix}
\PY{n}{p\PYZus{}reshaped} \PY{o}{=} \PY{n}{p}\PY{o}{.}\PY{n}{reshape}\PY{p}{(}\PY{o}{\PYZhy{}}\PY{l+m+mi}{1}\PY{p}{,} \PY{l+m+mi}{1}\PY{p}{)}
\PY{n}{y\PYZus{}reshaped} \PY{o}{=} \PY{n}{y}\PY{o}{.}\PY{n}{reshape}\PY{p}{(}\PY{o}{\PYZhy{}}\PY{l+m+mi}{1}\PY{p}{,} \PY{l+m+mi}{1}\PY{p}{)}

\PY{c+c1}{\PYZsh{} The score for the logit model is simply (y \PYZhy{} p)X}
\PY{n}{g} \PY{o}{=} \PY{p}{(}\PY{n}{y\PYZus{}reshaped} \PY{o}{\PYZhy{}} \PY{n}{p\PYZus{}reshaped}\PY{p}{)} \PY{o}{*} \PY{n}{X}

\PY{c+c1}{\PYZsh{} Calculate the BHHH estimator of the Information Matrix}
\PY{c+c1}{\PYZsh{} J = sum(g\PYZus{}i\PYZsq{} * g\PYZus{}i) }
\PY{n}{J} \PY{o}{=} \PY{n}{g}\PY{o}{.}\PY{n}{T} \PY{o}{@} \PY{n}{g}

\PY{c+c1}{\PYZsh{} The asymptotic variance\PYZhy{}covariance matrix is the inverse of J}
\PY{k}{try}\PY{p}{:}
    \PY{n}{asymptotic\PYZus{}var\PYZus{}cov\PYZus{}bhhh} \PY{o}{=} \PY{n}{np}\PY{o}{.}\PY{n}{linalg}\PY{o}{.}\PY{n}{inv}\PY{p}{(}\PY{n}{J}\PY{p}{)}

    \PY{c+c1}{\PYZsh{} The diagonal elements are the variances of the estimators}
    \PY{n}{variances\PYZus{}bhhh} \PY{o}{=} \PY{n}{np}\PY{o}{.}\PY{n}{diag}\PY{p}{(}\PY{n}{asymptotic\PYZus{}var\PYZus{}cov\PYZus{}bhhh}\PY{p}{)}

    \PY{c+c1}{\PYZsh{} The standard errors are the square roots of the variances}
    \PY{n}{std\PYZus{}errors\PYZus{}bhhh} \PY{o}{=} \PY{n}{np}\PY{o}{.}\PY{n}{sqrt}\PY{p}{(}\PY{n}{variances\PYZus{}bhhh}\PY{p}{)}

    \PY{n+nb}{print}\PY{p}{(}\PY{l+s+s2}{\PYZdq{}}\PY{l+s+s2}{Asymptotic Variance\PYZhy{}Covariance Matrix (BHHH estimator for Logit):}\PY{l+s+s2}{\PYZdq{}}\PY{p}{)}
    \PY{n+nb}{print}\PY{p}{(}\PY{n}{asymptotic\PYZus{}var\PYZus{}cov\PYZus{}bhhh}\PY{p}{)}
    \PY{n+nb}{print}\PY{p}{(}\PY{l+s+s2}{\PYZdq{}}\PY{l+s+se}{\PYZbs{}n}\PY{l+s+s2}{\PYZdq{}}\PY{p}{)}

    \PY{n+nb}{print}\PY{p}{(}\PY{l+s+s2}{\PYZdq{}}\PY{l+s+s2}{Asymptotic Variances of estimators (BHHH for Logit):}\PY{l+s+s2}{\PYZdq{}}\PY{p}{)}
    \PY{n+nb}{print}\PY{p}{(}\PY{l+s+sa}{f}\PY{l+s+s2}{\PYZdq{}}\PY{l+s+s2}{Var(beta\PYZus{}1\PYZus{}hat): }\PY{l+s+si}{\PYZob{}}\PY{n}{variances\PYZus{}bhhh}\PY{p}{[}\PY{l+m+mi}{0}\PY{p}{]}\PY{l+s+si}{:}\PY{l+s+s2}{.4f}\PY{l+s+si}{\PYZcb{}}\PY{l+s+s2}{\PYZdq{}}\PY{p}{)}
    \PY{n+nb}{print}\PY{p}{(}\PY{l+s+sa}{f}\PY{l+s+s2}{\PYZdq{}}\PY{l+s+s2}{Var(beta\PYZus{}2\PYZus{}hat): }\PY{l+s+si}{\PYZob{}}\PY{n}{variances\PYZus{}bhhh}\PY{p}{[}\PY{l+m+mi}{1}\PY{p}{]}\PY{l+s+si}{:}\PY{l+s+s2}{.4f}\PY{l+s+si}{\PYZcb{}}\PY{l+s+s2}{\PYZdq{}}\PY{p}{)}
    \PY{n+nb}{print}\PY{p}{(}\PY{l+s+s2}{\PYZdq{}}\PY{l+s+se}{\PYZbs{}n}\PY{l+s+s2}{\PYZdq{}}\PY{p}{)}

    \PY{n+nb}{print}\PY{p}{(}\PY{l+s+s2}{\PYZdq{}}\PY{l+s+s2}{Standard Errors of estimators (BHHH for Logit):}\PY{l+s+s2}{\PYZdq{}}\PY{p}{)}
    \PY{n+nb}{print}\PY{p}{(}\PY{l+s+sa}{f}\PY{l+s+s2}{\PYZdq{}}\PY{l+s+s2}{SE(beta\PYZus{}1\PYZus{}hat): }\PY{l+s+si}{\PYZob{}}\PY{n}{std\PYZus{}errors\PYZus{}bhhh}\PY{p}{[}\PY{l+m+mi}{0}\PY{p}{]}\PY{l+s+si}{:}\PY{l+s+s2}{.4f}\PY{l+s+si}{\PYZcb{}}\PY{l+s+s2}{\PYZdq{}}\PY{p}{)}
    \PY{n+nb}{print}\PY{p}{(}\PY{l+s+sa}{f}\PY{l+s+s2}{\PYZdq{}}\PY{l+s+s2}{SE(beta\PYZus{}2\PYZus{}hat): }\PY{l+s+si}{\PYZob{}}\PY{n}{std\PYZus{}errors\PYZus{}bhhh}\PY{p}{[}\PY{l+m+mi}{1}\PY{p}{]}\PY{l+s+si}{:}\PY{l+s+s2}{.4f}\PY{l+s+si}{\PYZcb{}}\PY{l+s+s2}{\PYZdq{}}\PY{p}{)}

\PY{k}{except} \PY{n}{np}\PY{o}{.}\PY{n}{linalg}\PY{o}{.}\PY{n}{LinAlgError}\PY{p}{:}
    \PY{n+nb}{print}\PY{p}{(}\PY{l+s+s2}{\PYZdq{}}\PY{l+s+s2}{Could not compute the variance\PYZhy{}covariance matrix.}\PY{l+s+s2}{\PYZdq{}}\PY{p}{)}
    \PY{n+nb}{print}\PY{p}{(}\PY{l+s+s2}{\PYZdq{}}\PY{l+s+s2}{The Information Matrix (J) may be singular.}\PY{l+s+s2}{\PYZdq{}}\PY{p}{)}
\end{Verbatim}
\end{tcolorbox}

    \begin{Verbatim}[commandchars=\\\{\}]
Asymptotic Variance-Covariance Matrix (BHHH estimator for Logit):
[[ 0.01624169 -0.01624169]
 [-0.01624169  0.02733954]]


Asymptotic Variances of estimators (BHHH for Logit):
Var(beta\_1\_hat): 0.0162
Var(beta\_2\_hat): 0.0273


Standard Errors of estimators (BHHH for Logit):
SE(beta\_1\_hat): 0.1274
SE(beta\_2\_hat): 0.1653
    \end{Verbatim}

    \subsubsection{Formulas for Asymptotic Variance (Logit
Model)}\label{formulas-for-asymptotic-variance-logit-model}

To calculate the asymptotic variance of the MLE estimators, we used the
\textbf{Berndt--Hall--Hall--Hausman (BHHH)} method. Here are the
formulas involved:

\paragraph{1. Score Vector (Gradient of the
Log-Likelihood)}\label{score-vector-gradient-of-the-log-likelihood}

First, for each observation \texttt{i}, we compute the score vector
\texttt{g\_i}, which is the gradient of the log-likelihood function for
that observation, evaluated at the MLE estimates
(\texttt{\textbackslash{}hat\{\textbackslash{}beta\}}).

The formula for the score \texttt{gᵢ} in the \textbf{Logit} model is
remarkably simple: \[ g_i = (y_i - \Lambda(X_i'\hat{\beta})) X_i \]

Where: - \texttt{yᵢ} and \texttt{Xᵢ} are the data for the i-th
observation. - \texttt{Λ} (Lambda) is the Cumulative Distribution
Function (CDF) of the standard logistic distribution.

\paragraph{2. Information Matrix (BHHH
Estimator)}\label{information-matrix-bhhh-estimator}

Next, we estimate the Information Matrix (\texttt{J}) by summing the
outer product of each observation's score vector with itself:

\[ J = \sum_{i=1}^{n} g_i g_i' \]

This is what the line \texttt{J\ =\ g.T\ @\ g} in the code calculates.

\paragraph{3. Asymptotic Variance-Covariance
Matrix}\label{asymptotic-variance-covariance-matrix}

Finally, the asymptotic variance-covariance matrix (\texttt{V}) is
calculated by taking the inverse of the Information Matrix \texttt{J}:

\[ V(\hat{\beta}) = J^{-1} \]

The diagonal elements of this resulting matrix \texttt{V} are the
variances for each of your estimated parameters
(\texttt{\textbackslash{}hat\{\textbackslash{}beta\}₁} and
\texttt{\textbackslash{}hat\{\textbackslash{}beta\}₂}). The square root
of these diagonal elements gives the standard errors.

    \begin{tcolorbox}[breakable, size=fbox, boxrule=1pt, pad at break*=1mm,colback=cellbackground, colframe=cellborder]
\prompt{In}{incolor}{5}{\boxspacing}
\begin{Verbatim}[commandchars=\\\{\}]
\PY{c+c1}{\PYZsh{} This cell calculates the 95\PYZpc{} confidence interval for beta\PYZus{}2\PYZus{}hat.}
\PY{c+c1}{\PYZsh{} It depends on \PYZsq{}mle\PYZus{}beta\PYZsq{} and \PYZsq{}std\PYZus{}errors\PYZus{}bhhh\PYZsq{} from the previous cells.}

\PY{c+c1}{\PYZsh{} The z\PYZhy{}score for a 95\PYZpc{} confidence interval is approximately 1.96}
\PY{n}{z\PYZus{}score} \PY{o}{=} \PY{l+m+mf}{1.96}

\PY{c+c1}{\PYZsh{} Get the MLE estimate and standard error for beta\PYZus{}2 from previous calculations}
\PY{n}{beta\PYZus{}2\PYZus{}hat} \PY{o}{=} \PY{n}{mle\PYZus{}beta}\PY{p}{[}\PY{l+m+mi}{1}\PY{p}{]}
\PY{n}{se\PYZus{}beta\PYZus{}2\PYZus{}hat} \PY{o}{=} \PY{n}{std\PYZus{}errors\PYZus{}bhhh}\PY{p}{[}\PY{l+m+mi}{1}\PY{p}{]}

\PY{c+c1}{\PYZsh{} Calculate the lower and upper bounds of the confidence interval}
\PY{n}{lower\PYZus{}bound} \PY{o}{=} \PY{n}{beta\PYZus{}2\PYZus{}hat} \PY{o}{\PYZhy{}} \PY{n}{z\PYZus{}score} \PY{o}{*} \PY{n}{se\PYZus{}beta\PYZus{}2\PYZus{}hat}
\PY{n}{upper\PYZus{}bound} \PY{o}{=} \PY{n}{beta\PYZus{}2\PYZus{}hat} \PY{o}{+} \PY{n}{z\PYZus{}score} \PY{o}{*} \PY{n}{se\PYZus{}beta\PYZus{}2\PYZus{}hat}

\PY{n+nb}{print}\PY{p}{(}\PY{l+s+s2}{\PYZdq{}}\PY{l+s+s2}{95}\PY{l+s+s2}{\PYZpc{}}\PY{l+s+s2}{ Confidence Interval for beta\PYZus{}2\PYZus{}hat (from Logit model):}\PY{l+s+s2}{\PYZdq{}}\PY{p}{)}
\PY{n+nb}{print}\PY{p}{(}\PY{l+s+sa}{f}\PY{l+s+s2}{\PYZdq{}}\PY{l+s+s2}{[}\PY{l+s+si}{\PYZob{}}\PY{n}{lower\PYZus{}bound}\PY{l+s+si}{:}\PY{l+s+s2}{.4f}\PY{l+s+si}{\PYZcb{}}\PY{l+s+s2}{, }\PY{l+s+si}{\PYZob{}}\PY{n}{upper\PYZus{}bound}\PY{l+s+si}{:}\PY{l+s+s2}{.4f}\PY{l+s+si}{\PYZcb{}}\PY{l+s+s2}{]}\PY{l+s+s2}{\PYZdq{}}\PY{p}{)}
\end{Verbatim}
\end{tcolorbox}

    \begin{Verbatim}[commandchars=\\\{\}]
95\% Confidence Interval for beta\_2\_hat (from Logit model):
[0.5654, 1.2136]
    \end{Verbatim}

    \begin{tcolorbox}[breakable, size=fbox, boxrule=1pt, pad at break*=1mm,colback=cellbackground, colframe=cellborder]
\prompt{In}{incolor}{6}{\boxspacing}
\begin{Verbatim}[commandchars=\\\{\}]
\PY{c+c1}{\PYZsh{} This cell runs a Monte Carlo simulation using the CORRECT Logit estimation.}
\PY{k+kn}{import}\PY{+w}{ }\PY{n+nn}{numpy}\PY{+w}{ }\PY{k}{as}\PY{+w}{ }\PY{n+nn}{np}
\PY{k+kn}{import}\PY{+w}{ }\PY{n+nn}{scipy}\PY{n+nn}{.}\PY{n+nn}{optimize}\PY{+w}{ }\PY{k}{as}\PY{+w}{ }\PY{n+nn}{opt}
\PY{k+kn}{from}\PY{+w}{ }\PY{n+nn}{scipy}\PY{+w}{ }\PY{k+kn}{import} \PY{n}{special}
\PY{k+kn}{from}\PY{+w}{ }\PY{n+nn}{tqdm}\PY{n+nn}{.}\PY{n+nn}{notebook}\PY{+w}{ }\PY{k+kn}{import} \PY{n}{tqdm} \PY{c+c1}{\PYZsh{} Optional: for a progress bar}

\PY{n}{n\PYZus{}simulations} \PY{o}{=} \PY{l+m+mi}{1000}
\PY{n}{beta\PYZus{}hats} \PY{o}{=} \PY{n}{np}\PY{o}{.}\PY{n}{zeros}\PY{p}{(}\PY{p}{(}\PY{n}{n\PYZus{}simulations}\PY{p}{,} \PY{l+m+mi}{2}\PY{p}{)}\PY{p}{)}
\PY{n}{beta\PYZus{}stes} \PY{o}{=} \PY{n}{np}\PY{o}{.}\PY{n}{zeros}\PY{p}{(}\PY{p}{(}\PY{n}{n\PYZus{}simulations}\PY{p}{,} \PY{l+m+mi}{2}\PY{p}{)}\PY{p}{)}

\PY{c+c1}{\PYZsh{} True parameter values from the problem description}
\PY{n}{beta\PYZus{}0} \PY{o}{=} \PY{n}{np}\PY{o}{.}\PY{n}{array}\PY{p}{(}\PY{p}{[}\PY{l+m+mi}{1}\PY{p}{,} \PY{l+m+mi}{1}\PY{p}{]}\PY{p}{)}

\PY{n+nb}{print}\PY{p}{(}\PY{l+s+sa}{f}\PY{l+s+s2}{\PYZdq{}}\PY{l+s+s2}{Running }\PY{l+s+si}{\PYZob{}}\PY{n}{n\PYZus{}simulations}\PY{l+s+si}{\PYZcb{}}\PY{l+s+s2}{ Monte Carlo simulations with correct Logit estimation...}\PY{l+s+s2}{\PYZdq{}}\PY{p}{)}

\PY{c+c1}{\PYZsh{} Loop for the Monte Carlo simulation}
\PY{k}{for} \PY{n}{i} \PY{o+ow}{in} \PY{n}{tqdm}\PY{p}{(}\PY{n+nb}{range}\PY{p}{(}\PY{n}{n\PYZus{}simulations}\PY{p}{)}\PY{p}{)}\PY{p}{:}
    \PY{c+c1}{\PYZsh{} 1. Draw a new sample}
    \PY{n}{y\PYZus{}draws}\PY{p}{,} \PY{n}{x\PYZus{}draws} \PY{o}{=} \PY{n}{LogitDraws}\PY{p}{(}\PY{p}{)}
    \PY{n}{X} \PY{o}{=} \PY{n}{np}\PY{o}{.}\PY{n}{hstack}\PY{p}{(}\PY{p}{(}\PY{n}{np}\PY{o}{.}\PY{n}{ones}\PY{p}{(}\PY{p}{(}\PY{n}{x\PYZus{}draws}\PY{o}{.}\PY{n}{shape}\PY{p}{[}\PY{l+m+mi}{0}\PY{p}{]}\PY{p}{,} \PY{l+m+mi}{1}\PY{p}{)}\PY{p}{)}\PY{p}{,} \PY{n}{x\PYZus{}draws}\PY{p}{)}\PY{p}{)}
    \PY{n}{y} \PY{o}{=} \PY{n}{y\PYZus{}draws}\PY{o}{.}\PY{n}{flatten}\PY{p}{(}\PY{p}{)}
    
    \PY{c+c1}{\PYZsh{} 2. Calculate MLE for beta using Logit likelihood}
    \PY{n}{initial\PYZus{}beta} \PY{o}{=} \PY{n}{np}\PY{o}{.}\PY{n}{zeros}\PY{p}{(}\PY{n}{X}\PY{o}{.}\PY{n}{shape}\PY{p}{[}\PY{l+m+mi}{1}\PY{p}{]}\PY{p}{)}
    \PY{n}{result} \PY{o}{=} \PY{n}{opt}\PY{o}{.}\PY{n}{minimize}\PY{p}{(}
        \PY{n}{neg\PYZus{}log\PYZus{}likelihood\PYZus{}logit}\PY{p}{,} \PY{c+c1}{\PYZsh{} Use the correct logit function}
        \PY{n}{initial\PYZus{}beta}\PY{p}{,}
        \PY{n}{args}\PY{o}{=}\PY{p}{(}\PY{n}{y}\PY{p}{,} \PY{n}{X}\PY{p}{)}\PY{p}{,}
        \PY{n}{method}\PY{o}{=}\PY{l+s+s1}{\PYZsq{}}\PY{l+s+s1}{BFGS}\PY{l+s+s1}{\PYZsq{}}
    \PY{p}{)}
    
    \PY{k}{if} \PY{o+ow}{not} \PY{n}{result}\PY{o}{.}\PY{n}{success}\PY{p}{:}
        \PY{n}{beta\PYZus{}hats}\PY{p}{[}\PY{n}{i}\PY{p}{,} \PY{p}{:}\PY{p}{]} \PY{o}{=} \PY{n}{np}\PY{o}{.}\PY{n}{nan}
        \PY{n}{beta\PYZus{}stes}\PY{p}{[}\PY{n}{i}\PY{p}{,} \PY{p}{:}\PY{p}{]} \PY{o}{=} \PY{n}{np}\PY{o}{.}\PY{n}{nan}
        \PY{k}{continue}
        
    \PY{n}{mle\PYZus{}beta} \PY{o}{=} \PY{n}{result}\PY{o}{.}\PY{n}{x}
    \PY{n}{beta\PYZus{}hats}\PY{p}{[}\PY{n}{i}\PY{p}{,} \PY{p}{:}\PY{p}{]} \PY{o}{=} \PY{n}{mle\PYZus{}beta}
    
    \PY{c+c1}{\PYZsh{} 3. Calculate Asymptotic Variance (Standard Errors) for Logit}
    \PY{n}{z} \PY{o}{=} \PY{n}{X} \PY{o}{@} \PY{n}{mle\PYZus{}beta}
    \PY{n}{p} \PY{o}{=} \PY{n}{special}\PY{o}{.}\PY{n}{expit}\PY{p}{(}\PY{n}{z}\PY{p}{)}
    \PY{n}{epsilon} \PY{o}{=} \PY{l+m+mf}{1e\PYZhy{}9}
    \PY{n}{p} \PY{o}{=} \PY{n}{np}\PY{o}{.}\PY{n}{clip}\PY{p}{(}\PY{n}{p}\PY{p}{,} \PY{n}{epsilon}\PY{p}{,} \PY{l+m+mi}{1} \PY{o}{\PYZhy{}} \PY{n}{epsilon}\PY{p}{)}
    
    \PY{n}{p\PYZus{}reshaped} \PY{o}{=} \PY{n}{p}\PY{o}{.}\PY{n}{reshape}\PY{p}{(}\PY{o}{\PYZhy{}}\PY{l+m+mi}{1}\PY{p}{,} \PY{l+m+mi}{1}\PY{p}{)}
    \PY{n}{y\PYZus{}reshaped} \PY{o}{=} \PY{n}{y}\PY{o}{.}\PY{n}{reshape}\PY{p}{(}\PY{o}{\PYZhy{}}\PY{l+m+mi}{1}\PY{p}{,} \PY{l+m+mi}{1}\PY{p}{)}
    
    \PY{c+c1}{\PYZsh{} Use the correct score for the Logit model: g = (y \PYZhy{} p)X}
    \PY{n}{g} \PY{o}{=} \PY{p}{(}\PY{n}{y\PYZus{}reshaped} \PY{o}{\PYZhy{}} \PY{n}{p\PYZus{}reshaped}\PY{p}{)} \PY{o}{*} \PY{n}{X}
    \PY{n}{J} \PY{o}{=} \PY{n}{g}\PY{o}{.}\PY{n}{T} \PY{o}{@} \PY{n}{g}
    
    \PY{k}{try}\PY{p}{:}
        \PY{n}{asymptotic\PYZus{}var\PYZus{}cov} \PY{o}{=} \PY{n}{np}\PY{o}{.}\PY{n}{linalg}\PY{o}{.}\PY{n}{inv}\PY{p}{(}\PY{n}{J}\PY{p}{)}
        \PY{n}{std\PYZus{}errors} \PY{o}{=} \PY{n}{np}\PY{o}{.}\PY{n}{sqrt}\PY{p}{(}\PY{n}{np}\PY{o}{.}\PY{n}{diag}\PY{p}{(}\PY{n}{asymptotic\PYZus{}var\PYZus{}cov}\PY{p}{)}\PY{p}{)}
        \PY{n}{beta\PYZus{}stes}\PY{p}{[}\PY{n}{i}\PY{p}{,} \PY{p}{:}\PY{p}{]} \PY{o}{=} \PY{n}{std\PYZus{}errors}
    \PY{k}{except} \PY{n}{np}\PY{o}{.}\PY{n}{linalg}\PY{o}{.}\PY{n}{LinAlgError}\PY{p}{:}
        \PY{n}{beta\PYZus{}stes}\PY{p}{[}\PY{n}{i}\PY{p}{,} \PY{p}{:}\PY{p}{]} \PY{o}{=} \PY{n}{np}\PY{o}{.}\PY{n}{nan}

\PY{n+nb}{print}\PY{p}{(}\PY{l+s+s2}{\PYZdq{}}\PY{l+s+s2}{Simulations finished.}\PY{l+s+s2}{\PYZdq{}}\PY{p}{)}

\PY{c+c1}{\PYZsh{} 4. Report how many times beta\PYZus{}0 was in the 95\PYZpc{} CI}
\PY{n}{z\PYZus{}score} \PY{o}{=} \PY{l+m+mf}{1.96}

\PY{c+c1}{\PYZsh{} Calculate lower and upper bounds for both parameters across all simulations}
\PY{n}{lower\PYZus{}bounds} \PY{o}{=} \PY{n}{beta\PYZus{}hats} \PY{o}{\PYZhy{}} \PY{n}{z\PYZus{}score} \PY{o}{*} \PY{n}{beta\PYZus{}stes}
\PY{n}{upper\PYZus{}bounds} \PY{o}{=} \PY{n}{beta\PYZus{}hats} \PY{o}{+} \PY{n}{z\PYZus{}score} \PY{o}{*} \PY{n}{beta\PYZus{}stes}

\PY{c+c1}{\PYZsh{} Check coverage for beta\PYZus{}1 (the intercept)}
\PY{n}{coverage\PYZus{}beta1} \PY{o}{=} \PY{n}{np}\PY{o}{.}\PY{n}{sum}\PY{p}{(}\PY{p}{(}\PY{n}{lower\PYZus{}bounds}\PY{p}{[}\PY{p}{:}\PY{p}{,} \PY{l+m+mi}{0}\PY{p}{]} \PY{o}{\PYZlt{}}\PY{o}{=} \PY{n}{beta\PYZus{}0}\PY{p}{[}\PY{l+m+mi}{0}\PY{p}{]}\PY{p}{)} \PY{o}{\PYZam{}} \PY{p}{(}\PY{n}{upper\PYZus{}bounds}\PY{p}{[}\PY{p}{:}\PY{p}{,} \PY{l+m+mi}{0}\PY{p}{]} \PY{o}{\PYZgt{}}\PY{o}{=} \PY{n}{beta\PYZus{}0}\PY{p}{[}\PY{l+m+mi}{0}\PY{p}{]}\PY{p}{)}\PY{p}{)}

\PY{c+c1}{\PYZsh{} Check coverage for beta\PYZus{}2 (the coefficient on X)}
\PY{n}{coverage\PYZus{}beta2} \PY{o}{=} \PY{n}{np}\PY{o}{.}\PY{n}{sum}\PY{p}{(}\PY{p}{(}\PY{n}{lower\PYZus{}bounds}\PY{p}{[}\PY{p}{:}\PY{p}{,} \PY{l+m+mi}{1}\PY{p}{]} \PY{o}{\PYZlt{}}\PY{o}{=} \PY{n}{beta\PYZus{}0}\PY{p}{[}\PY{l+m+mi}{1}\PY{p}{]}\PY{p}{)} \PY{o}{\PYZam{}} \PY{p}{(}\PY{n}{upper\PYZus{}bounds}\PY{p}{[}\PY{p}{:}\PY{p}{,} \PY{l+m+mi}{1}\PY{p}{]} \PY{o}{\PYZgt{}}\PY{o}{=} \PY{n}{beta\PYZus{}0}\PY{p}{[}\PY{l+m+mi}{1}\PY{p}{]}\PY{p}{)}\PY{p}{)}

\PY{n+nb}{print}\PY{p}{(}\PY{l+s+s2}{\PYZdq{}}\PY{l+s+se}{\PYZbs{}n}\PY{l+s+s2}{\PYZhy{}\PYZhy{}\PYZhy{} Coverage Results (Correct Logit Model) \PYZhy{}\PYZhy{}\PYZhy{}}\PY{l+s+s2}{\PYZdq{}}\PY{p}{)}
\PY{n+nb}{print}\PY{p}{(}\PY{l+s+sa}{f}\PY{l+s+s2}{\PYZdq{}}\PY{l+s+s2}{Out of }\PY{l+s+si}{\PYZob{}}\PY{n}{n\PYZus{}simulations}\PY{l+s+si}{\PYZcb{}}\PY{l+s+s2}{ iterations:}\PY{l+s+s2}{\PYZdq{}}\PY{p}{)}
\PY{n+nb}{print}\PY{p}{(}\PY{l+s+sa}{f}\PY{l+s+s2}{\PYZdq{}}\PY{l+s+s2}{The 95\PYZpc{} CI for the intercept (beta\PYZus{}1) contained the true value of }\PY{l+s+si}{\PYZob{}}\PY{n}{beta\PYZus{}0}\PY{p}{[}\PY{l+m+mi}{0}\PY{p}{]}\PY{l+s+si}{\PYZcb{}}\PY{l+s+s2}{ a total of }\PY{l+s+si}{\PYZob{}}\PY{n}{coverage\PYZus{}beta1}\PY{l+s+si}{\PYZcb{}}\PY{l+s+s2}{ times.}\PY{l+s+s2}{\PYZdq{}}\PY{p}{)}
\PY{n+nb}{print}\PY{p}{(}\PY{l+s+sa}{f}\PY{l+s+s2}{\PYZdq{}}\PY{l+s+s2}{The 95\PYZpc{} CI for the coefficient (beta\PYZus{}2) contained the true value of }\PY{l+s+si}{\PYZob{}}\PY{n}{beta\PYZus{}0}\PY{p}{[}\PY{l+m+mi}{1}\PY{p}{]}\PY{l+s+si}{\PYZcb{}}\PY{l+s+s2}{ a total of }\PY{l+s+si}{\PYZob{}}\PY{n}{coverage\PYZus{}beta2}\PY{l+s+si}{\PYZcb{}}\PY{l+s+s2}{ times.}\PY{l+s+s2}{\PYZdq{}}\PY{p}{)}
\PY{n+nb}{print}\PY{p}{(}\PY{l+s+sa}{f}\PY{l+s+s2}{\PYZdq{}}\PY{l+s+s2}{Expected coverage for a 95\PYZpc{} CI is approximately }\PY{l+s+si}{\PYZob{}}\PY{l+m+mf}{0.95}\PY{+w}{ }\PY{o}{*}\PY{+w}{ }\PY{n}{n\PYZus{}simulations}\PY{l+s+si}{\PYZcb{}}\PY{l+s+s2}{.}\PY{l+s+s2}{\PYZdq{}}\PY{p}{)}
\end{Verbatim}
\end{tcolorbox}

    \begin{Verbatim}[commandchars=\\\{\}]
Running 1000 Monte Carlo simulations with correct Logit estimation{\ldots}
    \end{Verbatim}

    
    \begin{Verbatim}[commandchars=\\\{\}]
  0\%|          | 0/1000 [00:00<?, ?it/s]
    \end{Verbatim}

    
    \begin{Verbatim}[commandchars=\\\{\}]
Simulations finished.

--- Coverage Results (Correct Logit Model) ---
Out of 1000 iterations:
The 95\% CI for the intercept (beta\_1) contained the true value of 1 a total of
955 times.
The 95\% CI for the coefficient (beta\_2) contained the true value of 1 a total of
964 times.
Expected coverage for a 95\% CI is approximately 950.0.
    \end{Verbatim}


    % Add a bibliography block to the postdoc
    
    
    
\end{document}
